\section{R ile hipotez testi}
\label{sec:r-b09}

\textbf{Soru.} Bölümdeki çift yönlü testte gözlenen ortalama farkı 2,1,
standart hata 1 ve serbestlik derecesi 49 olsun. $H_0:\delta=0$ için
\pval-değerini, yüzde 95 güven aralığını ve $\alpha=0{,}05$ kararını bulun.

\begin{lstlisting}[style=prismR]
fark <- 2.1
standart_hata <- 1
serbestlik <- 49
alfa <- 0.05
t_degeri <- fark / standart_hata
p_degeri <- 2 * pt(abs(t_degeri), df = serbestlik,
                   lower.tail = FALSE)
kritik <- qt(1 - alfa / 2, df = serbestlik)
aralik <- fark + c(-1, 1) * kritik * standart_hata
c(t = t_degeri, p = p_degeri, kritik = kritik)
aralik
if (p_degeri < alfa) {
  print("H0 reddedilir")
} else {
  print("H0 reddedilemez")
}
pt(t_degeri, df = serbestlik, lower.tail = FALSE)
\end{lstlisting}

\begin{outputguidebox}
\textbf{Çözüm ve kontrol değerleri.} $t=2{,}10$, çift yönlü
$p=0{,}04090$ ve kritik değer 2,0096'dır. Güven aralığı
$[0{,}0904;4{,}1096]$ sıfırı içermez; yüzde 5 düzeyinde $H_0$ reddedilir.
Son satır, önceden belirlenmiş $H_1:\delta>0$ için tek yönlü
$p=0{,}02045$ verir.
\end{outputguidebox}

\textbf{Yorum.} Çift yönlü değer, mutlak test istatistiğinin üst kuyruk
olasılığının iki katıdır. Sonuca bakıp yön seçmek uygun değildir. Küçük
\pval-değeri $H_0$'ın doğru olma olasılığı veya etkinin önemli olduğu anlamına gelmez.

\textbf{Pekiştirme ve yanıt.} Önceden $\alpha=0{,}01$ seçilseydi aynı
$p=0{,}04090$ ile $H_0$ reddedilemezdi. Veri değişmez; karar eşiği değişir.